Rotation-symmetric bosonic codes encode a qubit in the Fock states $\ket{mK}$ of a single oscillator. We follow a primitive that was found for an unrelated reason: the \textsc{spider} state, which maximally violates the classical bound of the precession protocol of Zaw \emph{et al.} [Phys.\ Rev.\ A \textbf{106}, 032222 (2022)] and happens to have exactly the required $K$-fold rotation symmetry. We show that the score operator of the protocol is block diagonal in the rotation-symmetry sectors, so that the spider state defines an order-$K$ code whose preparation can be certified device-independently from sign measurements alone, with a spectral-gap fidelity witness. The code, however, is not error-correcting: its codewords carry mean photon numbers in a fixed ratio of three, the Knill--Laflamme non-deformation condition fails by an amount of order the photon number, no finite reweighting of the coefficients can repair it without abandoning the maximal violation that defines the state, and the exact state has infinite mean energy. Since the obstacle to rotation-code error correction is then not the primitive but the gadgets, we turn to the controlled-rotation gate $e^{i\varphi\n\otimes\n}$ that every rotation code needs and engineer it from the drive-tuned cross-Kerr interaction between two superconducting cavities sharing a transmon. An untargeted drive enhances the cross-Kerr four-fold and yields a gate time of tens of microseconds, but produces self-Kerr of comparable size; a 99\% gate tolerates a residual self-Kerr below one percent of the cross-Kerr. We then reproduce, with the same dressed-ladder numerics, the single-tone drive-enhanced scheme that completed the project [Qin \emph{et al.}, arXiv:2609.07076]: a drive placed next to the two-cavity conversion resonance raises the cross-Kerr on-off ratio above fifty while the self-Kerrs stay small, and a driven side transmon---already present in recovered four-mode simulations---cancels the residual, leaving the data rail of the CROT gate exact in the diagonal model.
